\section{Calibration and Validation}\label{sec:calibration}

\subsection{Sector-specific neural calibration on the 31-ticker ladder}\label{sec:ladder}

The parametric $\psi$ of equation~\eqref{eq:psi} provided the analytic starting point but assumed a single global smile and term-structure shape across all tickers, an assumption that became a bottleneck on a heterogeneous universe spanning sectors with structurally different risk profiles. To quantify the gap and the value of relaxing it, we collected full option ladders for 31 tickers spanning six sectors on fifteen capture dates between 2026-04-14 and 2026-05-11, producing 234{,}549 observations after filtering to moneyness $K/S \in [0.80, 1.20]$, non-zero bids, and IV $\in (0.01, 2.0)$. The sector partition followed conventional GICS-style groupings (Supplementary Table~\ref{tab:sector_map}): technology (10 tickers), healthcare (8), financials (4), energy (3), retail (3), and broad-market ETFs (3). The corpus carried systematic cross-sector variation: mega-cap ETFs exhibited shallow smiles with pronounced term structure, technology names carried elevated near-term skew from earnings catalysts, healthcare names showed idiosyncratic bumps around trial and FDA dates, and commodity-linked energy tickers displayed skew patterns shaped by supply shocks rather than equity crash risk. Against this partition, we quantified the value of relaxing the global-shape assumption by fitting three nested models on the pooled ladder using a common loss (mean squared IV error) and a common level parameterization (one $\theta_{\text{ticker}}$ per ticker), varying only the representation of $\psi$: a global parametric $\psi$ with the five $\beta$-parameters shared across all tickers, a global neural $\psi_{\text{NN}}$ shared across all tickers that replaced the parametric form with the neural architecture introduced in the methods, and a sector-specific neural $\psi_{\text{NN}}^{(c)}$ with one network per sector. This design isolated the contribution of model capacity (parametric $\to$ NN) from the contribution of sharing granularity (global $\to$ sector).

We compared the three nested models on the pooled corpus and found that both axes contributed independently (Table~\ref{tab:three_way}): moving from the global parametric $\psi$ to the global neural $\psi$ reduced overall RMSE from 12.48\% to 11.47\% IV (an 8.1\% relative improvement) by accommodating smile shapes the log-polynomial form could not express, and moving from the global to the sector-specific $\psi_{\text{NN}}$ further reduced RMSE to 10.24\% (a 17.9\% cumulative relative improvement over the parametric baseline) by letting within-sector shapes specialize away from the global average. A per-day Hagan-2002 SABR reference fitted independently to each (date, ticker, maturity-bucket) slice (seven-day maturity buckets priced at their mean maturity, with the capture-day close as the forward proxy and no dividend adjustment) produced a pooled RMSE of 9.05\% IV across the same 234{,}549 observations; that figure undercuts the sector-NN $\psi$ by 1.2\% IV, which we read as the cost of cross-date pooling that lets $\psi$ serve forward simulation. SABR pays for its tighter pooled fit by refitting $\sim$6{,}072 parameters across the corpus (one $(\alpha,\rho,\nu)$ triple per slice with $\beta = 0.5$ fixed), and the resulting calibration is a snapshot of each (date, ticker, expiry) cross-section that cannot price an unseen date without re-fit, propagate forward under stochastic dynamics, or carry a shared shape across tickers in a sector; the sector $\psi_{\text{NN}}$ encodes a single calibrated surface (337 parameters per sector, six sectors) reusable across all 15 dates and at simulation time. These figures were in-sample fits to the pooled corpus; a leave-one-date-out check that held out the latest capture and retrained on the remaining fourteen days produced a pooled test RMSE of $10.89\%$ against a train RMSE of $10.21\%$, leaving a generalization gap of $+0.68\%$ IV (Supplementary Table~\ref{tab:loo_sector_nn}) and confirming that the sector-NN result reflected faithful generalization rather than overfit to the multi-date pool. We further audited the calibrated $\psi$ surfaces against the static no-arbitrage conditions $\partial^2 C/\partial K^2 \geq 0$ and $\partial C/\partial T \geq 0$ on a $20\times 20$ moneyness-maturity grid for six sector surfaces (one representative ticker per sector) and the two per-ticker GS and LLY surfaces used in the forward scenarios below; across the eight surfaces, 9.8\% of grid points (281/2880) violated the butterfly condition and 0.8\% (23/3040) violated the calendar condition, concentrated in the deep wings where ladder coverage was thinnest (Supplementary Table~\ref{tab:static_arb}). We report these rates as a known limitation of the unconstrained neural surrogate; introducing no-arbitrage penalties in the calibration loss is a natural extension reserved for future work.

\begin{table}[ht]
\centering
\caption{Overall RMSE comparison on the 31-ticker ladder (234{,}549 observations across fifteen capture dates). The parametric baseline fit a single global set of five $\beta$ parameters; the shared NN replaced the parametric $\psi$ with one global network; the sector NN used one network per sector. All three models fitted one $\theta_{\text{ticker}}$ per ticker jointly with the shape parameters. The per-day SABR reference fitted Hagan-2002 SABR ($\beta = 0.5$, free parameters $\alpha,\rho,\nu$) independently to every (date, ticker, maturity-bucket) slice and pooled the residuals (2{,}024 separate fits, $\sim$6{,}072 free parameters in aggregate); it shares no parameters across dates and provides no forward-generation mechanism, so its pooled RMSE is the practitioner reference against which the cross-date pooling cost of $\psi$ is measured rather than a peer model.}\label{tab:three_way}
\begin{tabular}{lrr}
\toprule
Model & RMSE (\% IV) & Improvement vs.\ parametric \\
\midrule
Parametric (5 $\beta$, global)     & 12.48 & (baseline) \\
Neural $\psi$ (1 network, global)  & 11.47 & 8.1\% \\
Neural $\psi$ (6 networks, sector) & 10.24 & 17.9\% \\
Per-day SABR (per slice)           & 9.05 & 27.5\% \\
\bottomrule
\end{tabular}
\end{table}

The per-sector decomposition exposed where each incremental model capacity helped most (Supplementary Table~\ref{tab:per_sector}). The global NN uniformly improved every sector over the parametric baseline, with the largest absolute gains in ETF (10.54\% $\to$ 8.69\%) and retail (11.73\% $\to$ 10.23\%), where the parametric log-polynomial struggled to fit shallow ETF smiles and idiosyncratic retail patterns simultaneously. Moving to sector-specific networks delivered a further 1 to 3\% IV reduction in ETF, financials, retail, and technology but barely moved the needle in healthcare (12.71 $\to$ 12.53\%) and energy (8.24 $\to$ 7.88\%), reflecting cross-ticker dispersion in those sectors that a single shared $\psi$ could not absorb even when freed from the parametric form. Technology remained the hardest fit overall (14.47\% with sector NN); inspection of per-ticker residuals (Supplementary Table~\ref{tab:worst_tickers}) traced this to heterogeneity within the sector itself rather than to a failure of the sector NN as a whole: within technology, names such as GOOG (9.57\%), AVGO (9.94\%), and NVDA (10.14\%) fit comfortably under the shared Tech surface, while QCOM (21.46\%), MU (18.45\%), INTC (17.78\%), and AMD (17.74\%) pulled the sector RMSE upward. The pattern was consistent with these outlier tickers carrying idiosyncratic smile shapes (elevated wings, asymmetric skew) that a shared Tech $\psi$ could not simultaneously accommodate without degrading the fit of the better-behaved names. The same story repeated in healthcare, where ABBV (9.13\%) and UNH (10.39\%) fit well but PFE (17.66\%), MRNA (17.49\%), and AMGN (13.36\%) drove the sector average. The per-ticker $\theta_{\text{ticker}}$ values recovered a sensible IV level ranking: high-variance names with concentrated risk (INTC, MRNA, MU, AMD) occupied the top of the sorted order at 94 to 117\% baseline IV, large-cap diversified names (AAPL, JPM, WMT) clustered near the middle at 46 to 52\%, and the broad-market ETFs (SPY, QQQ, IWM) sat at the low-variance end near 35 to 42\%. Within-sector clustering was evident in the sorted level estimates, confirming that sector membership captured a dominant axis of the cross-sectional IV structure.

Beyond the aggregate RMSE comparison, the residual structure under the sector NN was concentrated in the wings (moneyness below 0.90 and above 1.10) rather than near the money, consistent with the in-sample fits of the single-ticker baselines and with the known difficulty of pricing deep-OTM contracts where liquidity is thin and quoted IV is noisier. The learned $\psi$ surfaces were smooth in both $\ln\tau$ and $\ln m$ and displayed the expected qualitative features, including put skew ($\psi$ rising for moneyness below unity), smile curvature in the wings, and a mild term-structure gradient. The sector-to-sector differences in these surfaces motivated relaxing the global-shape assumption in the first place: the healthcare and technology surfaces curved more sharply in the near-term wings than the ETF surface, and the energy surface displayed a subtle asymmetry consistent with the commodity-linked risk profile of the sector. This calibration established the sector-specific neural $\psi$ as the operating point of the framework on the pooled in-sample corpus, but the multiplicative decomposition supported a finer relaxation: one $\psi_{\text{NN}}$ per ticker. We took up that extension on the same corpus below, where per-name sample sizes were large enough to fit per-ticker networks of the same architecture on nearly the full universe.

\subsection{Per-ticker neural calibration on the pooled corpus}\label{sec:per_ticker_nn}

We trained one $\psi_{\text{NN}}$ per qualified ticker on the full pooled corpus and recovered substantial within-sector shape detail that the sector network could not express (Fig.~\ref{fig:per_ticker_panels}). Twenty-nine of the thirty-one tickers cleared a 2{,}000-observation threshold sufficient to fit a $\psi_{\text{NN}}$ of the same architectural family as the sector level; the two unqualified tickers (PFE at 1{,}300 observations, BMY at 1{,}847) were held on the sector $\psi_{\text{NN}}$ as a fallback. Network width followed sample size: tickers with at least 5{,}000 observations received the sector-NN $2 \to 16 \to 16 \to 1$ architecture, smaller-$N$ qualified tickers used $2 \to 8 \to 8 \to 1$. The per-ticker $\psi_{\text{NN}}$ reduced the pooled in-sample RMSE from $10.24\%$ under the sector NN to $9.73\%$ under the combined model (a $5.0\%$ relative reduction) and improved twenty-one of twenty-nine qualified tickers, with the largest absolute gains on JNJ ($-3.97\%$), MRNA ($-3.76\%$), AAPL ($-1.71\%$), and INTC ($-1.49\%$)---all in sectors with the most heterogeneous within-sector smile geometry. The shape detail recovered by per-ticker fits (steeper put skew on NVDA than the shared Tech surface, gentler near-the-money curvature on MSFT, the upward right wing on LLY missing from the larger-cap-dominated Healthcare surface) is visible in Fig.~\ref{fig:per_ticker_panels}.

\begin{figure}[ht]
\centering
\includegraphics[width=\textwidth]{sections/figures/ladder_per_ticker_nn_smile_panels.pdf}
\caption{Per-ticker $\psi_{\mathrm{NN}}$ smile fits on the fifteen-date pooled corpus. Representative IV smile panels at mid-DTE for six qualified tickers spanning four sectors (SPY, NVDA, MSFT, LLY, GS, AVGO). The per-ticker network (solid black) tracked the smile shape more tightly than the sector network (dashed gray) on every panel, with the largest gap on NVDA and MSFT where the shared Tech surface flattened to accommodate AVGO's elevated wings.}\label{fig:per_ticker_panels}
\end{figure}

We translated the IV-RMSE story into a dollar-pricing story on the most recent capture (2026-05-11) by repricing each panel via 200-step CRR American at the IV implied by each calibration tier, and observed that the per-ticker $\psi_{\mathrm{NN}}$ collapsed dollar errors to within \$1 of market mid on liquid names while LLY carried a uniform IV-regime bias (Fig.~\ref{fig:price_error}). On SPY (DTE 7) and GS (DTE 11), the per-ticker $\psi_{\mathrm{NN}}$ sat inside the local bid-ask ribbon for $33\%$ and $88\%$ of contracts respectively, dominating both the parametric and sector tiers; on NVDA (DTE 7), all three tiers clustered within $\pm\$0.7$. LLY (DTE 11) showed a uniform positive bias of $+\$2$ to $+\$6$ across every model tier: the per-ticker $\psi_{\mathrm{NN}}$ encoded the time-averaged LLY IV across the fifteen capture dates, including the 04-28 drawdown day, so on a calmer 05-11 the static $\psi$ overestimated IV by several percentage points and the dollar error followed. The market-IV reference (yfinance IV repriced through CRR) sat at $\sim\$0.6$ on LLY, confirming the CRR pipeline was sound and the residual was an IV-regime mismatch rather than a pricing error---a failure mode the low-rank daily adapter subsequently addresses.

\begin{figure}[ht]
\centering
\includegraphics[width=\textwidth]{sections/figures/ladder_price_error_three_tickers.pdf}
\caption{Model-vs-market dollar pricing error (CRR American at the calibrated IV minus market mid) versus moneyness for SPY (DTE 7), NVDA (DTE 7), LLY (DTE 11), and GS (DTE 11) on the 2026-05-11 capture. Four model tiers are overlaid on each panel: parametric $\psi$ (light), sector $\psi_{\mathrm{NN}}$ (medium), per-ticker $\psi_{\mathrm{NN}}$ (dark), and market-IV reference (green dot-dash; yfinance IV repriced through CRR, near zero by construction). The gray ribbon was the local $\pm\tfrac{1}{2}$ bid--ask spread band, computed per panel by partitioning the contracts into twelve equal-width moneyness bins, taking the median half-spread $(\text{ask}-\text{bid})/2$ within each bin, and linearly interpolating the twelve bin values onto a fine moneyness grid. A point falling inside the band was a model error smaller than the typical bid--ask spread at that moneyness, i.e., a price within the range the market itself was quoting; the band widened toward the wings because bid--ask spreads were systematically larger for deep in- and out-of-the-money contracts. The per-ticker fit dominated parametric and sector tiers on SPY and GS, matched them on NVDA, and showed a uniform positive bias on LLY consistent with the per-ticker $\psi$ encoding the time-averaged IV across fifteen capture dates and the sampled day (05-11) being calmer than the calibration-window average.}\label{fig:price_error}
\end{figure}

Within the qualified set, eight tickers (BAC, JPM, WFC, CVX, OXY, XOM, UPS, WMT) did not gain over the sector NN; these were concentrated in financials, energy, and retail, sectors whose sector NN had already absorbed most of the within-sector smile dispersion. The temporal-generalization study below used an earlier eight-day window from the same capture (2026-04-14 through 04-24) with a fixed train/test split, so the holdout block was held apart from the pooled in-sample fits.

\subsection{Temporal holdout and earnings-aware calibration}\label{sec:earnings_holdout}

The pooled in-sample fit established that the sector-specific $\psi_{\text{NN}}$ could absorb the cross-sectional smile geometry of the 31-ticker universe, but it left the temporal-generalization question unanswered: would a $\psi$ trained on one block of capture days price contracts on a held-out subsequent block? To answer this, we used an eight-day subset of the pooled capture (2026-04-14 through 04-17 and 04-21 through 04-24, comprising 130{,}730 filtered observations) and split it into a six-day training block (04-14 through 04-22, 103{,}897 observations) and a two-day temporal holdout (04-23 and 04-24, 26{,}833 observations). Eight of the 31 tickers were absent from the first capture day (04-14 contained 23 names); every holdout ticker appeared on multiple training days, so the per-ticker $\theta_i$ learned on the training block transferred directly to the holdout. We then asked a sharper question: which fraction of any test-set generalization gap was attributable to scheduled corporate events embedded in the holdout window rather than to architectural overfit? To separate these contributions, we ran the sector $\psi_{\text{NN}}$ under three configurations: (A)~a pooled control identical to the in-sample architecture above (two inputs $\ln\tau,\,\ln m$, no exclusion); (B)~a non-earnings baseline that excluded all training and test rows for which the row's ticker or any same-sector peer was within $\pm 3$ days of an earnings print; and (C)~the earnings-aware model introduced in the methods that fed $\psi$ four inputs $(\ln\tau,\,\ln m,\,e,\,e^{\text{peer}})$ and used the same train and test rows as the control. The peer feature carried sector-coupling information directly: for any ticker $i$ on day $t$, $e^{\text{peer}}$ collapsed to zero if any same-sector equity printed earnings that day, regardless of $i$'s own distance to its next print.

\begin{table}[ht]
\centering
\caption{Temporal-holdout sector-NN comparison across four configurations on the same train (04-14..04-22, 103{,}897 obs) and test (04-23..04-24, 26{,}833 obs) split. The pooled control (A) reproduced the existing four-day reference number ($13.03\%$ test RMSE). Excluding earnings-window observations from both train and test (B) reduced the test RMSE to $7.96\%$ and collapsed the generalization gap from $+4.99$ to $-0.06\%$ IV. Adding the earnings indicator and peer-coupling feature to the input vector (C) on the original train and test rows reduced the pooled test RMSE to $11.98\%$ and the generalization gap to $+4.36\%$. The Horvath-style direct neural IV baseline (D) replaced the $\theta_i \cdot \psi$ decomposition with a single MLP whose input vector additionally carried a 31-ticker 8-dimensional embedding (1{,}753 total parameters); on the same train/test split it landed at $12.36\%$ test RMSE, intermediate between the 2-input pooled control and the 4-input earnings-aware sector NN.}\label{tab:earnings_holdout}
\resizebox{\textwidth}{!}{%
\begin{tabular}{lrrrrr}
\toprule
Configuration & $N_{\text{tr}}$ & $N_{\text{te}}$ & Train (\%) & Test (\%) & Gap (\%) \\
\midrule
A: pooled (2-input)         & 103{,}897 & 26{,}833 & 8.05 & 13.03 & $+4.99$ \\
B: non-earnings (2-input)   &  29{,}383 &  4{,}516 & 8.02 &  7.96 & $-0.06$ \\
C: earnings-aware (4-input) & 103{,}897 & 26{,}833 & 7.61 & 11.98 & $+4.36$ \\
D: Horvath direct (4-input + 8-d ticker embedding) & 103{,}897 & 26{,}833 & 8.29 & 12.36 & $+4.07$ \\
\bottomrule
\end{tabular}}
\end{table}

The three-way comparison resolved the generalization question (Table~\ref{tab:earnings_holdout}). Configuration~B tied the test-time generalization gap of the sector NN to scheduled earnings events: with all earnings-window observations removed from both splits, the test RMSE ($7.96\%$) fell to within rounding of the train RMSE ($8.02\%$), giving a generalization gap of $-0.06\%$ IV. The exclusion also shrank and recomposed both splits (train $103{,}897 \to 29{,}383$ rows, test $26{,}833 \to 4{,}516$), so the comparison bounds the event contribution rather than isolating it on matched rows; within that bound, the architecture generalized faithfully on non-event surfaces, and what failed was the model's ability to price options on days when an earnings catalyst dominated the IV surface. Configuration~C confirmed that adding event-aware inputs to the shape network recovered roughly one-fifth of the lost performance: relative to the pooled control, the test RMSE dropped from $13.03\%$ to $11.98\%$ ($1.05\%$ IV absolute, $8.1\%$ relative), and the generalization gap narrowed from $+4.99$ to $+4.36\%$. The peer-coupling feature did the heavy lifting in this reduction: the per-ticker breakdown showed the largest improvements on Tech names whose own earnings prints fell outside the immediate holdout window but whose sector peers printed during it (Supplementary Table~\ref{tab:tech_pivot}). QCOM (own print 6 days after the holdout) and AMD (12 days after) both benefited from the peer-coupling signal because INTC printed inside the window and dragged their IV surfaces with a $\sim 4\%$ IV reduction in test error each. The same mechanism reduced META, MSFT, and GOOG test RMSE by 2 to $4\%$ IV apiece. INTC itself, the print-day ticker, was the one Tech name whose RMSE rose under the earnings-aware fit ($33.17 \to 35.94\%$); its outsized $+2192\%$ EPS surprise produced an unanticipated post-print IV crush that no scheduled-event indicator could predict from the calendar alone.

A direct-neural baseline that replaced the $\theta_i \cdot \psi$ decomposition with a single MLP carrying a ticker embedding (Configuration~D, $12.36\%$ test RMSE) landed between the 2-input pooled control and the 4-input earnings-aware fit, confirming that the explicit per-ticker level prior in the $\theta_i \cdot \psi$ form contributes a measurable but modest inductive bias rather than the dominant share of the model's accuracy: a comparable-capacity direct net learned an effective level representation from the ticker embedding itself, and the residual gap to Configuration~C ($-0.38\%$ test RMSE) reflects the structural advantage of the level/shape factorization rather than a capacity ceiling. These results set the limit of what a calendar-driven event indicator can and cannot do. The limit was sharp: the indicator captured the part of event-day IV behaviour known in advance from the calendar (pre-earnings IV expansion on the reporting ticker, sympathy moves on same-sector peers), but it carried no information about the post-event direction or size of the IV move when the actual print deviated from consensus. The non-earnings baseline of Configuration~B was therefore the right reference point against which to compare future event mechanisms: any extension of the model aimed at pricing event-day IV surfaces had to be evaluated against the $7.96\%$ non-event RMSE rather than against the pooled $13.03\%$, because that gap quantified the residual event signal the current architecture could not reach. We return to this distinction in the discussion, where the failure of the indicator on INTC and its success on the sympathy peers motivate a per-ticker tail mechanism reserved for future work.
